\chapter{Теория}
\label{ch:intro}

\section{Дискретное по времени преобразование Фурье}

Для вычисления ПФ по дискретным отсчетам S(n) применяется ДВПФ

$$X(i\Omega n) = \sum_{n=-\infty}^{\infty}X(n)e^{-i\Omega n}$$

Здесь $\Omega$ - нормированная (цифровая) частота

$$\Omega = \frac{2\pi f}{f_s}[\frac{rad}{sample}]$$

Эта формула получается из формулы фазового сдвига между соседними дискретными отсчётами сигнала.

$$\omega T_s = 2\pi fT_s = \frac{2\pi f}{f_s}$$

В ДВПФ осуществляется переход от $\omega$ к $\Omega$ следующим образом

\begin{figure}[H]
    \centering
    \includegraphics[width=1.0\textwidth]{s1.png}
    \caption{Переход к цифровой частоте}
\end{figure} 

\section{Свойства ДВПФ}

\subsection{Периодичность}

ДВПФ имеет период $2\pi$

$$X(i\Omega + 2\pi k) = X(i\Omega)$$

\subsection{Комплексность}

ДВПФ можно записать как

$$X(i\Omega) = |X(i\Omega)|*e^{X(i\Omega)}$$

\subsection{Комплексно-сопряженность}

ДВПФ можно записать как

$$X(i\Omega) = X^*(-i\Omega)$$

\section{Реакция дискретной системы на дискретную комплексную экспоненту}

Подадим на вход системы $X(n) = e^{i\Omega_0n}$. Тогда на выходе получим

$$y(n) = \sum_{k}^{}X(k)h(n-k) = \sum_{k}^{}X(n-k)h(n) = \sum_{k}^{}h(k)e^{i\Omega_0(n-k)} = $$
$$\sum_{k}^{}h(k)e^{i\Omega_0n}e^{-i\Omega_0k} = e^{i\Omega_0n}\sum_{k}^{}h(k)e^{-i\Omega_0k} = $$
$$e^{i\Omega_0n}H(i\Omega_0)$$

$H(i\Omega_0)$ - комплексный коэффициент передачи цепи.Зная его, можно определить выходной сигнал, допустим

$$X(n) = Acos(\Omega_0n + \Theta)$$

Тогда сигнал на выходе

$$y(n) = |H(i\Omega_0)|*Acos(\Omega_0n + \Theta + \phi_0(\Omega_0))$$

На выходе получили сигнал той же формы, но с измененной амплитудой и частотой.

\section{Проблемы ДВПФ}


На практике для использования на ЭВМ ДВПФ не подходит, т.к 

\begin{enumerate}
    \item Требуется бесконечное кол-во отсчетов
    \item ДВПФ должно быть посчитано по непрерывной переменной $\Omega$
\end{enumerate}

Требуется определить ДВПФ в конечном числе точек частоты. Для этого ДВПФ разбивается на дискретные точки на интервале $(0,2\pi]$
на N точек с шагом $\Omega_k = K\frac{2\pi}{N} (k\in[0, N-1])$. Отсюда

$$X(\Omega_k) = \sum_{n=0}^{N-1}X(n)e^{-ik\frac{2\pi}{N}n}$$

$$\boxed{X(\Omega_k) = \sum_{n=0}^{N-1}X(n)e^{-\frac{-i2\pi nk}{N}}}$$

Расстояние между частотами - $\frac{2\pi}{N}$ \\

$$\boxed{f_k = k\frac{f_s}{N}}$$


\endinput